\chapter{Recursive dynamics}\label{chap:dynamics}

A dynamic run of \manref{cge_dynamic} solves the static model of
chapters~\ref{chap:equations} to~\ref{chap:calibration} for periods
$t = 0, 1, \dots, T$. Period 0 is the calibrated base. Agents do not look
ahead: each period is an equilibrium given the stocks and exogenous values of
that period, which are updated from the previous period's solution. The
updates follow the recursive dynamic extension of the standard model
\citep{thurlow2004}; the rates are chosen by the user, and none has a default
other than zero.

\section{Updates between periods}

Before period $t \ge 1$, in this order:

\begin{enumerate}
\item \textbf{Capital.} For each capital factor $f \in K$, with depreciation
rate $\delta_f$ and base gross rental rate $\rho_f$ (rental per unit of capital
at base prices, so that the base capital stock is $\V{YF}^0_f / \rho_f$), new
capital in the factor's units is
\begin{equation*}
N_{f,t} = \rho_f\,\omega_f\,\V{INV}_{t-1}, \qquad
\V{INV}_{t-1} = \sum_c \V{PQ}^0_c\,\V{QINV}_{c,t-1}, \qquad
\omega_f = \V{YF}^0_f \Big/ \sum_{f' \in K} \V{YF}^0_{f'}:
\end{equation*}
real investment at base prices, shared among capital factors by their base
income. If $f$ is specific to each activity,
\begin{equation*}
\V{QF}_{fa,t} = (1 - \delta_f)\,\V{QF}_{fa,t-1} + N_{f,t}\,\eta_{fa,t},
\end{equation*}
with allocation shares that favour activities whose rental rate is above the
average,
\begin{gather*}
\eta_{fa,t} \propto \max\Bigl[0,\; s_{fa}\Bigl(1 + \beta\Bigl(\frac{R_{fa}}{\bar R_f} - 1\Bigr)\Bigr)\Bigr],\\
s_{fa} = \frac{\V{QF}_{fa,t-1}}{\sum_{a'} \V{QF}_{fa',t-1}}, \qquad
R_{fa} = \V{W}_{f,t-1}\,\V{WFDIST}_{fa,t-1}, \qquad
\bar R_f = \sum_a s_{fa}\,R_{fa},
\end{gather*}
normalised to sum to one, where $\beta \ge 0$ is the sensitivity of investment
to rental differences (default 1). If $f$ is mobile,
$\V{QFS}_{f,t} = (1 - \delta_f)\,\V{QFS}_{f,t-1} + N_{f,t}$. A capital factor
may not be unemployed.
\item \textbf{Other factors.} $\V{QFS}_{f,t} = (1 + g_f)\,\V{QFS}_{f,t-1}$ for
a mobile factor, and $\V{QF}_{fa,t} = (1 + g_f)\,\V{QF}_{fa,t-1}$ for an
activity-specific one. The supply of an unemployed factor is endogenous, so
$g_f$ must be zero. The labour force of a factor in $FL$ grows at $g_f$
(chapter~\ref{chap:labour}).
\item \textbf{Population.} $\gamma_{ch,t} = (1 + g_P)\,\gamma_{ch,t-1}$ (and
$\gamma^h$ alike), $\V{HS}^0_t = (1 + g_P)\,\V{HS}^0_{t-1}$ and
$\V{POP}_t = (1 + g_P)\,\V{POP}_{t-1}$: subsistence consumption grows with the
population, and the health link compares health services per head. With a death rate, $g_P$ becomes $g_{P,t}$, and factors
that lose workers to deaths grow at $g_{f,t}$ (section~\ref{sec:population}).
\item \textbf{Government and foreign flows.} $\overline{\V{qg}}_c$,
$\V{trnsfr}_{d,gov}$, $\V{trnsfr}_{gov,gov}$ and $\V{trnsfr}_{shi,gov}$ grow at
$g_G$; flows in foreign currency ($\V{trnsfr}_{d,row}$, $\V{trnsfr}_{f,row}$,
$\V{trnsfr}_{gov,row}$, $\V{trnsfr}_{row,gov}$, $\V{trnsfr}_{shi,row}$, and
$\V{FSAV}$ when it is fixed) grow at $g_W$. Other fixed values, the numeraire
among them, are unchanged.
\item \textbf{Public debt}, when the model has debt
(section~\ref{sec:debt}). The deficit of period $t - 1$, $-\V{GSAV}_{t-1}$, is
borrowed, a share $\phi$ abroad:
\begin{equation*}
\V{DD}_t = \V{DD}_{t-1} + (1 - \phi)\,(-\V{GSAV}_{t-1}),
\qquad
\V{DF}_t = \V{DF}_{t-1} + \phi\,(-\V{GSAV}_{t-1}) / \V{EXR}_{t-1}.
\end{equation*}
A surplus repays debt in the same shares. Interest in period $t$ is paid on
these stocks (\eqref{eq:I3}, \eqref{eq:I9}, \eqref{eq:S3}).
\item \textbf{Shocks} for period $t$, if any, are applied by the user's
program; they may change the debt stocks and interest rates. A shock stays in
force in later periods unless the program changes it.
\item \textbf{Efficiency.} $\V{ad}_{a,t} = \lambda_{a,t}\,\V{ad}_{a,t-1}$, with
$\lambda_{a,t} = 1 + g^{ad}_a$; or a uniform $\lambda_t$ taken from a baseline
run; or, to calibrate a baseline to growth targets, the uniform $\lambda_t$ for
which real value added at base prices, $\sum_a \V{PVA}^0_a\,\V{QVA}_{a,t}$,
equals $(1 + g^Y_t)$ times its value in $t - 1$. That $\lambda_t$ is found by
secant iterations starting from $\lambda = 1$ and $\lambda = 1 + g^Y_t$, to a
relative gap of $10^{-13}$.
\end{enumerate}

Period $t$ is then solved from the solution of $t - 1$. Each period is stored
as an ordinary simulation, named with the run and the period, and tagged with
its run, period and efficiency factor.

\section{Baselines and scenarios}

A scenario that changes policy should reuse the baseline's efficiency factors
$\lambda_t$, so that efficiency does not respond to the scenario; its
results are then compared with the baseline period by period
(\manref{cge_dynresults}). Because the run is recursive, a static effect on
investment changes the capital stock of later periods: a policy that raises
consumption at the expense of investment can lower output in later periods,
which a static run does not show.

\section{Lags in user equations}

From version 0.9.2, user equations may refer to the values of variables in
earlier periods, \cmd{lag(VAR)} for $\V{VAR}_{t-1}$ and \cmd{lag(VAR, k)} for
$\V{VAR}_{t-k}$, $k \le 4$ (section~\ref{sec:userlags}). Before the first
period, the lagged values are those of the base. They extend the updates
between periods above with stocks and adjustment rules of the user's own.
